\chapter{Propuesta técnica}

\section{Sistema AMR-FOD-Kitting para línea HTP A320}

El sistema implanta robots móviles autónomos para suministro coordinado, trazabilidad e inspección preventiva. El AMR enlaza zonas de preparación, estaciones del HTP A320, áreas de retorno de útiles y puntos de supervisión. Su trabajo diario es simple de enunciar y exigente de ejecutar: llevar el elemento correcto al punto correcto y dejar evidencia digital de la misión.

\section{Modelo de robot y criterio de selección}

El diseño distingue cuatro niveles: modelo de simulación, modelo técnico seleccionado, benchmarks de mercado y plataforma final de compra. En CoppeliaSim se trabaja con una base tipo KUKA YouBot/KUKA Omnirob, suficiente para representar una cinemática omnidireccional con ruedas mecanum/suecas y ensayar maniobras laterales cerca de estaciones. Como referencia técnica de implantación se adopta una clase \textit{Robotnik RB-KAIROS/RB-KAIROS+}: AMR interior omnidireccional, arquitectura ROS 2 y carga útil de hasta 250 kg según la ficha técnica del fabricante \cite{robotnikRbKairos,robotnikRbKairosDatasheet}. AGILOX ODM se conserva como alternativa comercial de intralogística para dollies, con accionamiento omnidireccional, 300 kg de carga y velocidad máxima de 1,4 m/s \cite{agiloxOdm}.

MiR250 y Omron LD-250 no se presentan como la compra final, sino como benchmarks de AMR de 250 kg. El MiR250 publica 250 kg de carga, hasta 2 m/s, uso interior y superficie de carga compacta \cite{mir250Specs}; el Omron LD-250 publica 250 kg, 1,2 m/s, navegación natural, sensores de seguridad y E-stop \cite{omronLd250}. Esas referencias acotan coste, madurez y capacidades de mercado, pero no sustituyen el requisito de maniobra lateral de la plataforma objetivo.

\begin{longtable}{P{0.18\textwidth}P{0.25\textwidth}P{0.31\textwidth}P{0.16\textwidth}}
\toprule
\textbf{Rol en la propuesta} & \textbf{Modelo/referencia} & \textbf{Dato técnico relevante} & \textbf{Uso}\\
\midrule
\endhead
Simulación académica & KUKA YouBot / KUKA Omnirob o base mecanum equivalente & Ruedas mecanum/suecas, movimiento longitudinal, lateral y giro sobre el eje vertical. & Validar cinemática, rutas y aproximación a estación en CoppeliaSim.\\
Selección técnica & Robotnik RB-KAIROS / RB-KAIROS+ class & AMR interior omnidireccional, ROS 2, SLAM, 1,5 m/s, hasta 250 kg con ruedas mecanum fuertes, safety pack opcional \cite{robotnikRbKairosDatasheet}. & Referencia principal de plataforma objetivo.\\
Alternativa comercial & AGILOX ODM & AMR omnidireccional para dollies, 300 kg de carga y 1,4 m/s \cite{agiloxOdm}. & Opción de compra si se prioriza solución empaquetada de intralogística.\\
Benchmark comercial & MiR250 Base Robot & 250 kg de carga útil, hasta 2 m/s, uso interior y autonomía de turno según ficha del fabricante \cite{mir250Specs}. & Normalizar orden de magnitud de coste AMR.\\
Benchmark comercial & Omron LD-250 & 250 kg de carga útil, 1,2 m/s, SLAM, planificación de rutas, láseres de seguridad y E-stop \cite{omronLd250}. & Contrastar capacidad, seguridad y software de flota.\\
Referencia industrial pesada & KUKA omniMove & Plataforma omnidireccional para cargas XXL; KUKA cita transporte de piezas aeronáuticas como caso de uso \cite{kukaOmniMove}. & Justificar omnidireccionalidad en aeronáutica; se descarta por sobredimensionada.\\
\bottomrule
\caption{Trazabilidad entre modelo de simulación, referencias comerciales y plataforma objetivo.}
\label{tab:modelos-robot}
\end{longtable}

El robot se selecciona por la maniobra en estación, no por fabricación estructural. Debe entrar y salir sin giros amplios, mantenerse orientado hacia el punto de entrega, corregir lateralmente su posición y compartir pasillos con operarios. Esta decisión conecta con el temario: las ruedas omnidireccionales dan más maniobrabilidad que una base diferencial, aunque exigen mayor control, buena odometría, compensación de deslizamiento y validación de seguridad.

\section{Opciones de compra y decisión de ingeniería}

\begin{longtable}{P{0.22\textwidth}P{0.23\textwidth}P{0.25\textwidth}P{0.20\textwidth}}
\toprule
\textbf{Opción} & \textbf{Fortaleza} & \textbf{Riesgo/condición} & \textbf{Decisión}\\
\midrule
\endhead
RB-KAIROS-class & Coincide con cinemática omnidireccional, ROS 2, SLAM y carga de 250 kg. & Requiere cotización formal, safety pack y adaptación portakits. & \textbf{Base técnica recomendada}.\\
AGILOX ODM-class & Producto intralogístico maduro, omnidireccional y orientado a dollies. & Menos abierto para experimentación académica; integración depende de proveedor. & Alternativa comercial seria.\\
MiR250 / Omron LD-250 & Muy buenos benchmarks de AMR 250 kg y flota. & No resuelven igual el acoplamiento lateral si la estación exige orientación fija. & Referencia de coste/capacidad, no selección final.\\
KUKA omniMove & Referencia pesada y aeronáutica para cargas XXL. & Sobrecoste y sobredimensionamiento para kits, herramientas y consumibles. & No recomendado para este alcance.\\
\bottomrule
\caption{Opciones comerciales y decisión de ingeniería para la plataforma móvil.}
\label{tab:opciones-compra}
\end{longtable}

\section{Especificación de plataforma objetivo}

La plataforma objetivo se especifica por requisitos de operación, no por marca cerrada. El integrador podría seleccionar un chasis industrial equivalente si respeta la cinemática omnidireccional, la carga útil moderada y la separación entre navegación, seguridad y proceso. Así no se confunde el robot usado en simulación con el equipo que se compraría para una nave aeronáutica.

\begin{longtable}{P{0.22\textwidth}P{0.33\textwidth}P{0.35\textwidth}}
\toprule
\textbf{Decisión de diseño} & \textbf{Especificación propuesta} & \textbf{Justificación técnica}\\
\midrule
\endhead
Locomoción & Base holonómica con cuatro ruedas mecanum/suecas, control de \(v_x\), \(v_y\) y \(\omega\). & Permite aproximación lateral a estaciones, corrección fina sin reorientar la carga y menor ocupación de pasillo que una base diferencial.\\
Carga y módulo & Clase 200--300 kg con bandeja portakits, puntos de amarre, balizas y posibilidad de cajones o estantería baja. & Cubre kits, herramientas, consumibles y útiles moderados sin plantear manipulación estructural del HTP.\\
Motorización & Motores BLDC o servomotores eléctricos con reductora, encoder por rueda, supervisión de corriente y temperatura. & El temario de motores favorece accionamientos eléctricos compactos, eficientes y con realimentación para control de velocidad/posición.\\
Control de bajo nivel & Conversión cinemática inversa de velocidad de cuerpo a velocidad de rueda; lazo cerrado por encoder; límites de aceleración. & La omnidireccionalidad mejora la maniobra, pero penaliza si no se controla el deslizamiento y la saturación de motores.\\
Control de navegación & Fusión LiDAR--IMU--encoders para pose; SLAM/mapa; ruta global A*/Dijkstra; evitación local DWA o campos potenciales. & Conecta la arquitectura con localización, planificación y anticolisión del curso sin vender un algoritmo como garantía universal.\\
Seguridad funcional & Escáner láser de seguridad, bumpers, E-stop físico, campos de protección, reducción de velocidad por zona y estado seguro ante fallo. & La seguridad no depende del LiDAR de navegación ni de la cámara FOD; queda como capa independiente de parada.\\
Energía y disponibilidad & Batería industrial con retorno automático a carga, reserva mínima operativa y mantenimiento preventivo de ruedas/sensores. & Los actuadores consumen batería y se degradan; el plan de flota debe considerar autonomía, carga y desgaste.\\
Interfaz e integración & HMI/tablet, WiFi industrial, base de datos de misiones y conectores progresivos hacia MES/ERP. & El piloto debe dejar evidencia operativa y preparar el crecimiento sin exigir integración total desde el primer día.\\
\bottomrule
\caption{Especificación técnica de la plataforma AMR omnidireccional objetivo.}
\label{tab:especificacion-amr-objetivo}
\end{longtable}

La decisión técnica mantiene tres niveles: \textit{YouBot/Omnirob} para simular cinemática, MiR250/LD-250 para acotar el orden de magnitud comercial y AMR omnidireccional industrial como plataforma final. Con esa trazabilidad, el presupuesto no fuerza una compra incompatible con el criterio técnico. Las referencias de 250 kg dimensionan coste y carga; la licitación exigiría base holonómica o mecanum equivalente.

\section{Funciones principales}

\begin{itemize}
  \item Transporte interno de kits, herramientas, consumibles y útiles de carga moderada.
  \item Lectura RFID/QR de bandejas, herramientas, órdenes y puntos de entrega.
  \item Entrega confirmada mediante HMI, tablet de estación o escaneo de operario autorizado.
  \item Inspección preventiva FOD en rutas o zonas definidas.
  \item Retorno de útiles usados y aviso de incidencias de material.
  \item Registro de misiones, eventos, tiempos y excepciones para análisis posterior.
\end{itemize}

\section{Sensores y componentes}

\begin{longtable}{P{0.28\textwidth}P{0.62\textwidth}}
\toprule
\textbf{Elemento} & \textbf{Justificación}\\
\midrule
\endhead
LiDAR 2D/3D & Navegación, SLAM y detección geométrica de obstáculos. No se trata como único elemento de seguridad funcional.\\
Cámara RGB-D o industrial & Inspección FOD preventiva, verificación visual de zonas y evidencia fotográfica de incidencias, siempre con validación humana ante positivos.\\
RFID/QR & Trazabilidad de kits, herramientas, entregas, devoluciones y puntos de estación.\\
IMU y encoders & Odometría y fusión sensorial para estimar pose con menor deriva.\\
Bumpers y escáner láser de seguridad & Capa separada de seguridad funcional: zonas de protección, reducción de velocidad y parada segura ante presencia humana.\\
Parada de emergencia & Requisito básico de operación industrial y aceptación por seguridad.\\
HMI/tablet & Solicitud de misiones, confirmación de entregas y visualización de estado.\\
WiFi industrial & Comunicación con supervisor, gestor de flota y base de datos.\\
Base de datos de misiones & Registro de eventos, trazabilidad, KPIs, histórico FOD y auditoría operativa.\\
\bottomrule
\end{longtable}

\section{Arquitectura sensorial por capas}

La arquitectura separa sensores de navegación, seguridad y proceso. El LiDAR de navegación y la fusión encoders--IMU localizan el AMR y mantienen el mapa; el escáner láser de seguridad, bumpers y E-stop forman una capa independiente de parada segura; la cámara RGB-D, RFID y QR cubren FOD y trazabilidad. Una cámara FOD no certifica calidad aeronáutica por sí sola. Un LiDAR de navegación tampoco sustituye un sistema de seguridad funcional.

\begin{longtable}{P{0.20\textwidth}P{0.24\textwidth}P{0.22\textwidth}P{0.24\textwidth}}
\toprule
\textbf{Capa} & \textbf{Componentes} & \textbf{Variable medida} & \textbf{Decisión que habilita}\\
\midrule
\endhead
Localización & Encoders, IMU, LiDAR 2D/3D & Odometría, orientación, nube/distancia a obstáculos. & Estimar pose, actualizar mapa y ejecutar SLAM.\\
Planificación & Mapa, gestor de misiones, LiDAR & Ruta global, ruta local, zonas bloqueadas. & Usar A*/Dijkstra para ruta global y DWA/campos potenciales para evitación local.\\
Seguridad & Escáner láser de seguridad, bumpers, E-stop & Invasión de campo protegido, contacto, alarma manual. & Reducir velocidad, parar, bloquear misión o retirar AMR.\\
Proceso & RGB-D/industrial, RFID, QR, HMI & Evidencia visual, identificador de kit, punto de entrega y confirmación humana. & Registrar entrega, ronda FOD, excepción y trazabilidad.\\
Supervisión & WiFi industrial, base de datos, dashboard & Estado de flota, tiempos, alarmas, KPIs. & Priorizar misiones, auditar eventos y decidir escalado.\\
\bottomrule
\caption{Capas sensoriales y decisiones operativas asociadas.}
\label{tab:capas-sensores}
\end{longtable}

\section{Actuadores, motorización y control}

La plataforma se concibe como un AMR eléctrico de ruedas omnidireccionales, coherente con los contenidos del curso sobre robots que ruedan, efectores, actuadores y controladores. Cada rueda mecanum/sueca se acciona mediante motor eléctrico con controlador de velocidad, realimentación por encoder y supervisión de corriente/temperatura. El control de bajo nivel convierte consignas de velocidad longitudinal, lateral y angular en velocidades de rueda, mientras que el nivel de navegación limita aceleraciones, radios de seguridad y velocidades máximas por zona. Esta base se selecciona por maniobrabilidad, no por capacidad de carga extrema: la misión prevista es mover kits, herramientas y útiles moderados.

Los actuadores no están previstos para manipular piezas estructurales del HTP. Su función es mover la base, accionar elementos auxiliares simples del módulo portacargas si se incorporan, activar señales luminosas/acústicas y ejecutar parada segura. Esta separación entre locomoción, seguridad y proceso evita sobredimensionar el sistema y mantiene el alcance dentro de una aplicación acotada de robótica móvil industrial.

\section{Modo de operación seguro}

El AMR circula a velocidad limitada en zonas compartidas, reduce velocidad al aproximarse a estaciones, mantiene zonas de protección configurables y ejecuta parada segura si se invade el campo de seguridad. Las rutas se validan en CoppeliaSim y posteriormente en planta con pruebas controladas. La supervisión permite bloquear zonas, priorizar misiones urgentes y retirar un robot de servicio si aparece una alarma.
